% =============================================================================
% PLANTILLA CORPORATIVA - COMPUTAEX
% Archivo principal: main.tex
% =============================================================================
\documentclass[aspectratio=169]{beamer} % Permite cambiar la relación de aspecto
\usepackage{computaex_theme} 

% --- DATOS DE LA PRESENTACIÓN ---
\title[Título Corto]{Título Completo de la Presentación}
\subtitle{Subtítulo}
\author[Nombre Corto]{Nombre Completo Autor}
\institute[COMPUTAEX]{Fundación COMPUTAEX}
\date{\today}

\begin{document}

% --- PORTADA (Sin logo en la esquina para evitar duplicidad) ---
{
\setbeamertemplate{logo}{}
\begin{frame}
\titlepage
\end{frame}
}

% --- ÍNDICE ---
\begin{frame}
\frametitle{Contenido de la Sesión}
\tableofcontents % Genera el índice automáticamente basado en las \section
\end{frame}

% =============================================================================
\section{Introducción y Objetivos}
% =============================================================================

\begin{frame}
\frametitle{Diapositiva de Texto y Listas}
Este es un ejemplo de cómo organizar puntos clave en una diapositiva:
\begin{itemize}
    \item \textbf{Punto importante:} Descripción breve del concepto.
    \item \textbf{Punto secundario:} Detalle adicional que ayuda a la comprensión.
    \begin{itemize}
        \item Subpunto de nivel 2.
        \item Otro subpunto.
    \end{itemize}
    \item \textbf{Conclusión parcial:} Lorem ipsum.
\end{itemize}
\end{frame}

% =============================================================================
\section{Análisis de Datos y Gráficos}
% =============================================================================

\begin{frame}
\frametitle{Uso de Bloques y Columnas}
Dividir la pantalla para comparar información o añadir una imagen lateral.
\vspace{0.3cm}

\begin{columns}[T] % T para alinear el texto arriba
    \begin{column}{0.45\textwidth}
        \begin{block}{Resaltado A}
            Contenido destacado.
        \end{block}
        
        \begin{exampleblock}{Caso de Éxito}
            Ejemplo práctico de aplicación del concepto.
        \end{exampleblock}
    \end{column}
    
    \begin{column}{0.45\textwidth}
        \begin{alertblock}{Advertencia o Límite}
            Información crítica o restricciones técnicas.
        \end{alertblock}
        \textit{Nota: Puedes insertar una imagen aquí con \texttt{\textbackslash includegraphics}.}
    \end{column}
\end{columns}
\end{frame}

% ==========================================
\section{Presentación de Datos}
% ==========================================

\begin{frame}
\frametitle{Tablas Comparativas}
Para presentar datos técnicos, se recomienda el uso del paquete \texttt{booktabs} para evitar líneas verticales innecesarias y mejorar la lectura:

\vspace{0.3cm}
\begin{table}
\centering
\begin{tabular}{lrc}
\toprule
\textbf{Algoritmo} & \textbf{Consultas (n=1)} & \textbf{Complejidad} \\
\midrule
Clásico Determínico & 2 & $O(n)$ \\
Clásico Probabilista & 1 & $O(1)$ \\
\textbf{Cuántico (Deutsch)} & \textbf{1} & \textbf{$O(1)$} \\
\bottomrule
\end{tabular}
\caption{Comparativa de eficiencia en la resolución del problema.}
\end{table}
\end{frame}

% =============================================================================
\section{Ecuaciones y Código}
% =============================================================================

\begin{frame}
\frametitle{Matemáticas y Entornos Científicos}
Para fórmulas complejas o definiciones técnicas, usamos el modo matemático:

\vspace{0.2cm}
\begin{equation}
    E = mc^2 \quad \longrightarrow \quad |\psi\rangle = \alpha|0\rangle + \beta|1\rangle
\end{equation}

\vspace{0.4cm}
También puedes usar el entorno de descripción:
\begin{description}
    \item[Variable $x$:] Representa el valor de entrada del sensor.
    \item[Variable $y$:] Representa el flujo de datos procesado.
\end{description}
\end{frame}

% ==========================================
\section{Ejemplo Circuito Cuántico}
% ==========================================

\begin{frame}[fragile] % IMPORTANTE: [fragile] es obligatorio para usar quantikz en Beamer
\frametitle{Creación de Entrelazamiento: El Estado de Bell}
Para demostrar las capacidades gráficas de la plantilla, representamos el circuito de Bell:

\vspace{0.5cm}
\begin{center}
    \begin{quantikz}
        \lstick{$\ket{0}$} & \gate{H} & \ctrl{1} & \qw & \meter{} \\
        \lstick{$\ket{0}$} & \qw      & \targ{}  & \qw & \meter{}
    \end{quantikz}
\end{center}

\end{frame}

% ==========================================
\section{Cierre}
% ==========================================

\begin{frame}
\frametitle{Conclusiones y Siguientes Pasos}
Resumen final de la intervención:
\begin{enumerate}
    \item Primero, hemos analizado el estado actual del proyecto.
    \item Segundo, se han propuesto mejoras técnicas.
    \item Por último, establecemos el calendario de ejecución.
\end{enumerate}
\end{frame}

% ==========================================
\section{Bibliografía}
% ==========================================

\begin{frame}[allowframebreaks] % 'allowframebreaks' por si la lista es muy larga
\frametitle{Fuentes y Referencias}
\begin{thebibliography}{9}
    \setbeamercolor{bibliography item}{fg=MiColor}
    \setbeamercolor{bibliography entry author}{fg=black}

    \bibitem{nielsen_chuang}
    M. Nielsen \& I. Chuang.
    \textit{Quantum Computation and Quantum Information}.
    Cambridge University Press, 2010.

    \bibitem{computaex_rep}
    Don Pepito.
    \textit{Manual de ejemplos inventados}.
    Empresa ficticia, 2024.
\end{thebibliography}
\end{frame}

% --- DIAPOSITIVA FINAL ---
\begin{frame}
\begin{center}
    \Huge \textcolor{MiColor}{\textbf{Gracias por su atención}}
    
    \vspace{1cm}
    \large ¿Tienen alguna pregunta?
    
    \vspace{1.5cm}
    \small \textbf{Contacto:} \texttt{usuario@computaex.es} \\
    \textit{Fundación COMPUTAEX}
\end{center}
\end{frame}

\end{document}